% ECONOMICS_PROBLEM: {"id": "F13-420", "title": "Reshoring and friend-shoring welfare", "jel_code": "F13", "source_id": "OP-0088"}
% !TeX program = xelatex
\documentclass[11pt,a4paper]{article}
\usepackage[margin=23mm]{geometry}
\usepackage{fontspec}
\setmainfont{Latin Modern Roman}
\usepackage{amsmath,amssymb,mathtools}
\usepackage{xcolor,enumitem,booktabs,longtable,array,xurl,hyperref}
\definecolor{muted}{HTML}{53636C}
\hypersetup{colorlinks=true,urlcolor=blue,linkcolor=blue}
\setlength{\parindent}{0pt}
\setlength{\parskip}{5pt}
\newcommand{\R}{\mathbb R}
\newcommand{\E}{\mathbb E}
\newcommand{\N}{\mathbb N}
\newcommand{\ind}{\mathbf 1}
\newcommand{\Var}{\operatorname{Var}}
\newcommand{\Cov}{\operatorname{Cov}}
\newcommand{\supp}{\operatorname{supp}}
\DeclareMathOperator*{\argmin}{arg\,min}
\DeclareMathOperator*{\argmax}{arg\,max}
\newcommand{\entrymeta}[3]{{\sffamily\small\color{muted}\textbf{\texttt{#1}}\quad|\quad #2\par #3\par}}
\newcommand{\Problemref}[1]{\texttt{#1}}
\newcommand{\JELref}[2]{\texttt{#2} (JEL \texttt{#1})}
\begin{document}
\section*{Reshoring and friend-shoring welfare}
\textbf{Problem ID:} \texttt{F13-420}\quad \textbf{JEL:} \texttt{F13}\par
% BEGIN ECONOMICS STATEMENT
\label{problem:OP-0088}
\label{atlas:prob:T8}

\noindent\textbf{Atlas ID:} T8; \textbf{original number:} 88; \textbf{field:} International trade. \textbf{Classification:} Editorial research formulation (theoretical/empirical); not a certified unsolved theorem.

\subsection*{Definitions and assumptions}
A firm selects supplier weights $x\in\Delta^J=\{x\ge0:\sum_jx_j=1\}$, capacity and inventories before stochastic disruptions $Z$ are realized. Model $\theta$ specifies production substitution, investment costs, delivery failures and the distribution $\mathbb P$ of $Z$, including cross-location dependence. Consumers choose feasible consumption, and firms and households determine a selected market equilibrium under sourcing policy $b$. Let $W_\theta(b;\mathbb P)$ be finite expected social utility net of real resilience costs under explicit welfare weights. Tariffs or subsidies must have a stated financing and revenue-recycling rule.

\subsection*{Formal research question}
Let $\mathcal B$ be a compact, prespecified class of implementable reshoring or diversification rules and $\mathcal P$ a nonempty set of plausible disruption laws. Determine the robust policy problem
\[
 \sup_{b\in\mathcal B}\inf_{\mathbb P\in\mathcal P}
  [W_\theta(b;\mathbb P)-W_\theta(b^0;\mathbb P)],
\]
or its identified range when technological parameters are not point identified. Establish conditions for existence and characterize when geographic diversification, domestic redundancy or unconstrained sourcing attains the highest welfare. Estimate willingness to pay for reliability separately from market power, subsidies, insurance and firms’ subjective disaster beliefs.

\subsection*{Interpretation and scope}
Diversification across strongly correlated suppliers may provide little insurance; reshoring can also concentrate common domestic hazards. A security objective must enter as an explicit payoff or constraint rather than an unexplained welfare gain. Rare disruption probabilities are weakly identified by short panels, so stress-test results depend on the stated law or ambiguity set. Private revealed willingness to pay need not equal social value when shortages spill over to other firms or taxpayers provide support. Comparing policies requires endogenous prices, entry and foreign retaliation. Existence needs continuity, compactness and controlled expectations; these are conditions to prove for a proposed model, not automatic properties of resilience.

\subsection*{Source audit and status}
The probable source behind ``Grossman–Helpman (2021)'' is [S1], whose full authors also include Hugo Lhuillier; its published 2023 title adds ``International.'' The fragment remains only provisionally disambiguated. [S2] documents sourcing changes and [S3] verifies an EU diversification objective. The atlas’s ``production-network literature'' is not a bibliographic source. These foundations do not certify current open status.

\subsection*{References}

\begin{enumerate}[label={[S\arabic*]},leftmargin=*]

\item Gene M. Grossman, Elhanan Helpman, and Hugo Lhuillier (2021). \emph{Supply Chain Resilience: Should Policy Promote Diversification or Reshoring?}. NBER Working Paper 29330; published with International added to the title in Journal of Political Economy 131(12), 3462–3496 (2023). \url{https://www.nber.org/papers/w29330}.
\textit{Locator and role:} Publisher or author abstract / bibliographic record; Supply-chain resilience and welfare of diversification versus reshoring. Probable disambiguation: original citation omitted Hugo Lhuillier and a title. The author-year fragment alone does not prove that this was the intended source.

\item Laura Alfaro and Davin Chor (2023). \emph{Global Supply Chains: The Looming “Great Reallocation”}. NBER Working Paper 31661. \url{https://doi.org/10.3386/w31661}.
\textit{Locator and role:} Publisher or author abstract / bibliographic record; Changing sourcing patterns and limits of direct bilateral trade data. Provides background or a benchmark result; does not establish that the editorial question remains unresolved in 2026.

\item European Commission (2026). \emph{EU Trade Relations with China}. Directorate-General for Trade and Economic Security, institutional webpage. \url{https://policy.trade.ec.europa.eu/eu-trade-relationships-country-and-region/countries-and-regions/china_en}.
\textit{Locator and role:} Publisher or author abstract / bibliographic record; Stated policy of reducing critical dependencies and diversifying where necessary. Verifies the institutional objective only; supplies no welfare estimate or probability distribution for geopolitical disruption. Accessed 8 October 2026.

\item \textbf{Unresolved original citation:} ``production-network literature.'' No author, title or year was supplied, so no unique source can be assigned.

\end{enumerate}

\subsection*{Common conventions: Source mathematical conventions}

Unless an entry states otherwise, agent and firm sets are finite. State and action spaces are measurable, strategies and outcome maps are measurable, probability laws are specified, and expectations exist for the targets under discussion. Infinite-horizon discount factors lie in $(0,1)$ and discounted objectives are finite. Norms, tolerances and complexity input sizes must be specified for computational comparisons. Constants or primitives that appear locally are inputs, not empirical facts asserted by this catalogue.

\textbf{Identification.} A statistical model $m\in\mathcal M$ implies an observable law $P_m$ and target $\theta(m)$. At law $P$, its identified set is
\[
\Theta(P;\mathcal M)=\{\theta(m):m\in\mathcal M,\ P_m=P\}.
\]
Point identification holds when this set is a singleton, or equivalently when $P_{m_1}=P_{m_2}$ implies $\theta(m_1)=\theta(m_2)$ for all compatible models. Sharp bounds mean bounds attained, or approached when the boundary is not attained, by compatible models. Writing this set is a definition, not a solution: the substantive task is to establish informative restrictions and derive or estimate the set. An unrestricted-confounding model may give only trivial bounds.

\textbf{Inference.} Identification, estimability and valid uncertainty quantification are separate questions. Fix $\alpha\in(0,1)$. A confidence procedure $C_n$ over a declared law class $\mathcal P$ has asymptotic uniform coverage at level $1-\alpha$ when
\[
\liminf_{n\to\infty}\inf_{P\in\mathcal P}\Pr_P\{\theta(P)\in C_n\}\geq1-\alpha.
\]
For set-identified targets the coverage event must be defined separately, for example coverage of the full identified set. If an entry uses identification-indexed models instead of uniquely indexed laws, the same coverage convention is applied over those models. Pointwise limits do not establish uniform validity, and asymptotic validity does not guarantee finite-sample precision. Sample sizes, dependence structures and weak-identification sequences are part of each application.

\textbf{Counterfactuals.} $Y_i(a)$ is a potential outcome under a specified intervention, including its timing, implementation and versions. It is not $Y_i$ conditional on voluntarily choosing $a$. A full intervention vector permits interference; shorthand scalar treatment notation does not silently impose no interference. Assignment, selection, attrition, anticipation and measurement must be specified. A claim about an instrument's local effect is not automatically a population policy effect.

\textbf{Economic equilibrium.} A counterfactual model includes best responses, constraints, feasibility, market clearing, state transitions and expectations consistent with the stated information. When several equilibria exist, either a declared selector is an input or the target is set-valued. Comparisons across policy regimes must use the stated selection convention. Reduced-form relationships are not presumed invariant after an equilibrium-changing intervention.

\textbf{Optimization and welfare.} Optimization is over the stated feasible set. If attainment is not established, $\sup$ or $\inf$ and $\varepsilon$-optimal choices replace an asserted maximizer or minimizer. For example, with value $V=\sup_{a\in\mathcal A}W(a)<\infty$, an $\varepsilon$-optimal set is $\{a\in\mathcal A:W(a)\geq V-\varepsilon\}$ for $\varepsilon>0$. Benefits, costs, risk and health or environmental outcomes can be aggregated only using declared units and conversion weights. Interpersonal utility weights, the treatment of fiscal transfers and foreign profits, discounting, and the behavioral welfare criterion are explicit inputs. A comparative-static derivative additionally requires existence and appropriate differentiability of the selected counterfactual.

\textbf{Sources.} Local labels [S1], [S2], and so forth restart in every question. Locators identify the relevant abstract, model, section or result at the level actually checked; a bibliographic or abstract-level verification is not represented as inspection of an entire proof. Working papers are distinguished from later publications and changing versions. Source summaries are paraphrased. All URL checks and editorial formulations refer to the audit date stated above.
% END ECONOMICS STATEMENT
\end{document}
